\chapter{Pary uporządkowane i relacje}\label{ch:relations}

\section{Od nieuporządkowanych zbiorów do uporządkowanych par}
Korzystamy z pojęć i praw działań na zbiorach
z rozdziału~\ref{ch:sets}. Zbiór $\{a,b\}$ nie rozróżnia kolejności: zachodzi
$\{a,b\}=\{b,a\}$. Tymczasem punkt płaszczyzny z pierwszą
współrzędną $2$ i drugą $5$ powinien być inny niż punkt o
współrzędnych $5$ i $2$. Podobnie zapis „$a$ pozostaje w relacji z
$b$” może być prawdziwy bez prawdziwości zdania
„$b$ pozostaje w relacji z $a$”. Potrzebujemy zatem pojęcia pary,
w której pozycja elementu ma znaczenie.

\begin{definicja}[Para uporządkowana --- konstrukcja Kuratowskiego]
\label{def:kuratowski-pair}
Dla dowolnych zbiorów $a,b$ definiujemy
\[
(a,b):=\bigl\{\{a\},\{a,b\}\bigr\}.
\]
Jest to \emph{para uporządkowana} o pierwszej współrzędnej $a$
i drugiej współrzędnej $b$. Konstrukcję tę wybieramy jako
jedno z możliwych kodowań uporządkowanej pary za pomocą zbiorów.
\end{definicja}

\begin{twierdzenie}[Jednoznaczność współrzędnych pary]
\label{thm:kuratowski}
Dla dowolnych zbiorów $a,b,c,d$ zachodzi
\[
(a,b)=(c,d)\quad\Longleftrightarrow\quad a=c\text{ i }b=d.
\]
\end{twierdzenie}
\begin{proof}
Jeżeli $a=c$ i $b=d$, podstawienie do definicji
bezpośrednio daje równość par.

Załóżmy odwrotnie, że $(a,b)=(c,d)$. Zauważmy, że
\[
\bigcap\{\{a\},\{a,b\}\}=\{a\},\qquad
\bigcup\{\{a\},\{a,b\}\}=\{a,b\}.
\]
Są to przecięcie i suma elementów \emph{niepustej} rodziny
kodującej parę. Analogiczne równości zachodzą dla $c,d$.
Równość par oznacza równość kodujących je rodzin zbiorów,
a więc równość ich przecięć i sum. Otrzymujemy
$\{a\}=\{c\}$ oraz $\{a,b\}=\{c,d\}$.
Z pierwszej równości wynika $a=c$.

Aby wywnioskować $b=d$ z $\{a,b\}=\{a,d\}$,
rozważmy dwa przypadki. Jeśli $b=a$, to lewa strona
jest singletonem $\{a\}$. Równość wymusza
$d\in\{a\}$, zatem $d=a=b$. Jeżeli $b\ne a$,
to $b\in\{a,d\}$ i $b\ne a$, więc $b=d$.
W obu przypadkach druga współrzędna jest jednoznaczna.
\end{proof}

\begin{przyklad}[Dwa przypadki kodowania pary]\label{ex:kuratowski-cases-xlii}
Jeżeli $a=b$, to $\{a,b\}=\{a\}$, zatem
$(a,a)=\{\{a\}\}$, czyli kod jest zbiorem
jednoelementowym, którego jedynym elementem jest $\{a\}$.
Jeżeli $a\ne b$, to $\{a\}\ne\{a,b\}$ i kod
$(a,b)=\{\{a\},\{a,b\}\}$ ma dwa różne elementy.
Nie należy mylić liczby elementów \emph{samego kodu}
z liczbą współrzędnych: uporządkowana para zawsze
ma dwie współrzędne, nawet jeśli są one równe.
\end{przyklad}


\begin{uwaga}[Dlaczego definicja jest umowna?]
W wielu działach matematyki nie korzystamy z wnętrza kodowania
$(a,b)$; potrzebujemy wyłącznie własności
z Twierdzenia~\ref{thm:kuratowski}. Można wybrać inne
kodowanie spełniające to samo kryterium. Takie rozróżnienie
między obiektem a jego reprezentacją będzie ważne również
przy budowie liczb naturalnych w teorii ZF.
\end{uwaga}

\begin{definicja}[Iloczyn kartezjański]\label{def:cartesian}
Dla zbiorów $A,B$ definiujemy
$A\times B:=\{(a,b):a\in A,\ b\in B\}$.
Jeśli $A=\varnothing$ lub $B=\varnothing$, to
$A\times B=\varnothing$.
Dla $n\ge2$ możemy analogicznie rozważać uporządkowane
$n$-tki: definiujemy je rekurencyjnie, np.
$(a_1,a_2,a_3):=((a_1,a_2),a_3)$, a dla większego
$n$ dodajemy kolejne współrzędne. Iloczyn
$A_1\times\cdots\times A_n$ składa się z takich
$n$-tek, dla których $a_i\in A_i$ dla każdego $i$.
Warto zapamiętać, że zapis trójki jest tutaj kodowaniem,
a nie utożsamieniem dwóch różnych nawiasowań.
\end{definicja}

\begin{twierdzenie}[Dlaczego iloczyn kartezjański jest zbiorem?]
\label{thm:cartesian-existence-xlii}
Dla dowolnych zbiorów $A,B$ oraz $U:=A\cup B$ zachodzi
\[
A\times B\subseteq\calP\bigl(\calP(U)\bigr).
\]
W konsekwencji, gdy dostępne są operacje tworzenia pary,
sumy, zbioru potęgowego i wycinania podzbioru według
warunku, zbiór $A\times B$ rzeczywiście istnieje.
\end{twierdzenie}
\begin{proof}
Weźmy dowolną parę $(a,b)\in A\times B$. Z definicji
$a\in A$ i $b\in B$, a więc oba należą do $U$.
Wobec tego $\{a\}\subseteq U$ i $\{a,b\}\subseteq U$.
Definicja zbioru potęgowego daje
$\{a\}\in\calP(U)$ i $\{a,b\}\in\calP(U)$.
Konstrukcja Kuratowskiego mówi, że
$(a,b)=\{\{a\},\{a,b\}\}$; oba elementy tego zbioru
należą do $\calP(U)$, stąd $(a,b)\subseteq\calP(U)$.
Ponownie korzystając z definicji zbioru potęgowego,
otrzymujemy $(a,b)\in\calP(\calP(U))$.
Dowolność pary daje żądane zawieranie.

Mamy zatem \emph{ustalony zbiór kandydatów}
$\calP(\calP(U))$. Wycinając z niego te elementy $z$,
dla których istnieją $a\in A$ i $b\in B$ spełniające
$z=\{\{a\},\{a,b\}\}$, otrzymujemy dokładnie
$A\times B$. Zasada wycinania nie wymaga
nieograniczonego tworzenia zbiorów. W ZF jest to
zastosowanie odpowiednich aksjomatów, które
omówimy w rozdziale~\ref{ch:zf-axioms}.
\end{proof}


\begin{przyklad}
Jeśli $A=\{1,2\}$ i $B=\{x,y\}$, to
$A\times B=\{(1,x),(1,y),(2,x),(2,y)\}$.
Zwykle $A\times B\ne B\times A$, nawet gdy
$|A|=|B|$. Na przykład $(1,x)\in A\times B$,
ale $(1,x)\notin B\times A$ przy założeniu, że
$1,2,x,y$ są parami różne.
\end{przyklad}

\begin{twierdzenie}[Iloczyn zachowuje zawieranie]
\label{thm:cart-monotonic}
Niech $A,A',B,B'$ będą zbiorami. Jeżeli
$A\subseteq A'$ i $B\subseteq B'$, to
$A\times B\subseteq A'\times B'$.
Ponadto dla dowolnych zbiorów $A,B,C$ zachodzi
\[
A\times(B\cup C)=(A\times B)\cup(A\times C),\qquad
A\times(B\cap C)=(A\times B)\cap(A\times C).
\]
\end{twierdzenie}
\begin{proof}
Weźmy $(a,b)\in A\times B$. Wtedy $a\in A$, $b\in B$,
a więc z założeń $a\in A'$ oraz $b\in B'$.
Zatem $(a,b)\in A'\times B'$, co dowodzi zawierania.

Dla sumy ustalmy dowolną parę $(a,b)$.
Należy ona do $A\times(B\cup C)$ wtedy i tylko wtedy,
gdy $a\in A$ oraz $b\in B\cup C$, co znaczy,
że $a\in A$ i ($b\in B$ lub $b\in C$).
To dokładnie warunek $(a,b)\in A\times B$ lub
$(a,b)\in A\times C$. Dla przecięcia
zamieniamy w tym samym rozumowaniu „lub” na „i”.
Kryterium jednoznaczności par
(Twierdzenie~\ref{thm:kuratowski}) gwarantuje, że każda para
w iloczynie jest rozpoznawana po swoich współrzędnych.
\end{proof}

\begin{uwaga}[Wielokrotne iloczyny]
Istnieje kanoniczna bijekcja między $(A\times B)\times C$
a $A\times(B\times C)$, wysyłająca $((a,b),c)$ na
$(a,(b,c))$, ale przy przyjętym kodowaniu uporządkowanych par
nie muszą to być dosłownie równe zbiory.
Zwyczajowe zapisy trójek ukrywają tę różnicę.
\end{uwaga}

\begin{przyklad}[Ile jest par w skończonym iloczynie?]
Jeśli $A=\{a_1,a_2\}$ oraz $B=\{b_1,b_2,b_3\}$,
to dla każdego z dwóch pierwszych elementów mamy trzy
możliwe drugie współrzędne; otrzymujemy sześć par.
Ogólnie, jeżeli $A$ ma $m$ elementów, a $B$ ma
$n$ elementów, to $A\times B$ ma $mn$ elementów.
Równość nie wymaga, aby $A$ i $B$ były rozłączne,
gdyż pozycja współrzędnej rozstrzyga o parze.
\end{przyklad}

\section{Definicja relacji}\label{sec:relation-definition}
\begin{definicja}[Relacja binarna]\label{def:binary-relation}
Niech $A,B$ będą zbiorami. \emph{Relacją binarną z $A$ do $B$}
nazywamy dowolny podzbiór $R\subseteq A\times B$.
Zapis $aRb$ jest skrótem warunku $(a,b)\in R$.
Jeśli $A=B$, mówimy również o relacji \emph{na $A$}.
Dla ustalonej dodatniej liczby naturalnej $n$ relacją
$n$-argumentową na zbiorach $A_1,\ldots,A_n$ nazywamy
dowolny podzbiór $A_1\times\cdots\times A_n$;
na potrzeby tego zapisu jednoelementową krotkę
rozumiemy po prostu jako jej jedyną współrzędną.
\end{definicja}

\begin{przyklad}[Liczba możliwych relacji]\label{ex:number-relations-xlii}
Niech $A=\{1,2\}$ i $B=\{u,v\}$, gdzie wszystkie
wskazane elementy w odpowiednich zbiorach są różne.
Iloczyn $A\times B$ składa się z czterech par:
$(1,u),(1,v),(2,u),(2,v)$. Relacja z $A$ do $B$
jest \emph{dowolnym podzbiorem} tego iloczynu.
Dla każdej z czterech par niezależnie podejmujemy
decyzję „należy” albo „nie należy”, co daje $2^4=16$
różnych relacji. Wśród nich są relacja pusta, pełna
i na przykład relacja zawierająca tylko $(1,u)$.
Ogólnie dla skończonych $A,B$ o odpowiednio $m,n$
elementach istnieje dokładnie $2^{mn}$ relacji
z $A$ do $B$: wynika to z Twierdzenia~\ref{thm:powerset-finite}
i liczby elementów iloczynu kartezjańskiego.
Nie każda z tych relacji jest funkcją.
\end{przyklad}



\begin{przyklad}[Mniejsze niż]
Zbiór $R_<:=\{(m,n)\in\Z\times\Z:m<n\}$
jest relacją na $\Z$. Zapis $mR_<n$ jest po prostu
innym zapisem $m<n$. W szczególności $2R_<3$,
ale $3R_<2$ nie zachodzi.
\end{przyklad}

\begin{przyklad}[Podzielność i wybór zbioru bazowego]
Na $\N_{>0}:=\{1,2,3,\ldots\}$ relacja
$m\mid n\iff \exists k\in\N_{>0}\colon n=mk$
ma inne własności niż podzielność na $\Z$.
Na przykład $2\mid 6$ i $3\mid 6$, ale $2\nmid3$.
Na $\Z$ mamy $2\mid(-2)$ i $(-2)\mid2$ przy $2\ne-2$.
Będzie to ważny kontrprzykład do antysymetrii.
\end{przyklad}

\begin{definicja}[Dziedzina i zbiór wartości relacji]
\label{def:relation-dom-ran}
Niech $A,B$ będą zbiorami i $R\subseteq A\times B$. Określamy
\[
\dom(R)=\{a\in A:\exists b\in B\ (a,b)\in R\},\qquad
\ran(R)=\{b\in B:\exists a\in A\ (a,b)\in R\}.
\]
Zbiór $\dom(R)$ zawiera argumenty występujące na pierwszej,
a $\ran(R)$ na drugiej współrzędnej jakiejś pary z $R$.
\end{definicja}

\begin{przyklad}[Relacja z luką w dziedzinie]
Niech $A=\{1,2,3\}$, $B=\{u,v\}$ i
$R=\{(1,u),(1,v),(3,v)\}$.
Wtedy $\dom(R)=\{1,3\}$ oraz $\ran(R)=\{u,v\}$.
Element $2$ należy do ustalonego $A$, ale nie do dziedziny
relacji $R$. Co więcej, jednemu elementowi $1$ przypisano
\emph{dwa} elementy $u,v$. Relacja nie musi więc być funkcją;
precyzyjną definicję funkcji podamy w rozdziale o odwzorowaniach.
\end{przyklad}

\begin{definicja}[Obraz i przeciwobraz względem relacji]
\label{def:relation-image}
Niech $A,B$ będą zbiorami, $R\subseteq A\times B$,
$X\subseteq A$ i $Y\subseteq B$. Definiujemy
\[
R[X]:=\{b\in B:\exists a\in X\ (a,b)\in R\},\qquad
R^{-1}[Y]:=\{a\in A:\exists b\in Y\ (a,b)\in R\}.
\]
Nazywamy je odpowiednio \emph{obrazem $X$} i
\emph{przeciwobrazem $Y$ względem relacji}.
Dla ogólnej relacji przeciwobraz definiujemy więc przez
\emph{istnienie} powiązanego elementu. Gdy relacja będzie wykresem
funkcji, odzyskamy zwykłe przeciwobrazy funkcji. W szczególności
nie wolno automatycznie przenosić wszystkich wzorów znanych
z przeciwobrazów funkcji na dowolne relacje.
\end{definicja}

\begin{twierdzenie}[Obrazy relacji i działania na zbiorach]
\label{thm:relation-image-unions}
Niech $A,B,I$ będą zbiorami, $R\subseteq A\times B$,
a $(X_i)_{i\in I}$ rodziną podzbiorów $A$. Wtedy
\[
R\left[\bigcup_{i\in I}X_i\right]
=\bigcup_{i\in I}R[X_i].
\]
Dla $X,Z\subseteq A$ zachodzi ponadto
$R[X\cap Z]\subseteq R[X]\cap R[Z]$.
\end{twierdzenie}
\begin{proof}
Dla $b\in B$ przynależność do lewej strony pierwszej
równości oznacza istnienie $a\in\bigcup_i X_i$
takiego, że $(a,b)\in R$. Z definicji sumy istnieje
zatem indeks $i$ z $a\in X_i$. Otrzymujemy $b\in R[X_i]$,
a więc $b$ należy do prawej strony. Odwrotnie,
jeżeli $b\in R[X_i]$ dla pewnego $i$, istnieje
$a\in X_i\subseteq\bigcup_jX_j$ z $(a,b)\in R$;
stąd $b$ należy do lewej strony.

Jeżeli $b\in R[X\cap Z]$, wybierzmy $a\in X\cap Z$
z $(a,b)\in R$. Ten sam element $a$ świadczy jednocześnie
o $b\in R[X]$ i $b\in R[Z]$. Stąd drugie zawieranie.
\end{proof}

\begin{definicja}[Jednoznaczność po obu stronach relacji]
\label{def:relation-unique-fibers-xlii}
Niech $A,B$ będą zbiorami i niech $R\subseteq A\times B$. Relację $R$ nazwiemy
\emph{prawostronnie jednoznaczną}, gdy dla wszystkich
$a\in A$ i $b_1,b_2\in B$ warunki $aRb_1$ i $aRb_2$
implikują $b_1=b_2$. Relację $R$ nazwiemy
\emph{lewostronnie jednoznaczną}, gdy dla wszystkich
$b\in B$ i $a_1,a_2\in A$ warunki $a_1Rb$ i $a_2Rb$
implikują $a_1=a_2$. Nie żądamy tu istnienia
żadnej pary dla każdego elementu $A$ lub $B$.
\end{definicja}

\begin{twierdzenie}[Kiedy relacyjne obrazy zachowują przecięcia?]
\label{thm:relation-images-intersections-xlii}
Niech $A,B$ będą dowolnymi zbiorami, a
$R\subseteq A\times B$. Następujące warunki są
równoważne:
\begin{enumerate}[label=(\roman*)]
\item $R$ jest lewostronnie jednoznaczna;
\item dla wszystkich $X,Z\subseteq A$ zachodzi
$R[X\cap Z]=R[X]\cap R[Z]$.
\end{enumerate}
Ponadto następujące warunki są równoważne:
\begin{enumerate}[label=(\roman*)]
\item $R$ jest prawostronnie jednoznaczna;
\item dla wszystkich $Y,W\subseteq B$ zachodzi
$R^{-1}[Y\cap W]=R^{-1}[Y]\cap R^{-1}[W]$.
\end{enumerate}
\end{twierdzenie}
\begin{proof}
Załóżmy najpierw lewostronną jednoznaczność. Znamy już
zawieranie $R[X\cap Z]\subseteq R[X]\cap R[Z]$
(Twierdzenie~\ref{thm:relation-image-unions}).
Dla zawierania odwrotnego weźmy $b\in R[X]\cap R[Z]$.
Istnieje $a_1\in X$ takie, że $a_1Rb$, i istnieje
$a_2\in Z$ takie, że $a_2Rb$. Lewostronna jednoznaczność
daje $a_1=a_2$. Ten wspólny element należy więc do
$X\cap Z$ i jest związany z $b$, stąd $b\in R[X\cap Z]$.
To dowodzi równości.

Załóżmy teraz równość obrazów przecięć dla \emph{wszystkich}
$X,Z\subseteq A$. Przypuśćmy, że $a_1Rb$ i $a_2Rb$.
Gdyby $a_1\ne a_2$, to dla singletonów
$X=\{a_1\}$ i $Z=\{a_2\}$ mielibyśmy
$X\cap Z=\varnothing$, a zatem $R[X\cap Z]=\varnothing$.
Z drugiej strony $b$ należy do obu zbiorów $R[X]$
i $R[Z]$, więc jest elementem ich przecięcia,
sprzecznie z założoną równością. Koniecznie $a_1=a_2$.

Dla drugiej części dowód prowadzimy z pełną kontrolą
orientacji par. Przy prawostronnej jednoznaczności
weźmy $a\in R^{-1}[Y]\cap R^{-1}[W]$.
Istnieje $b_1\in Y$ z $aRb_1$ i istnieje $b_2\in W$
z $aRb_2$. Jednoznaczność daje $b_1=b_2$,
więc ten element leży w $Y\cap W$, skąd
$a\in R^{-1}[Y\cap W]$. Przeciwne zawieranie
wynika bez założeń z definicji obrazu relacji odwrotnej.

Wreszcie jeśli druga równość obowiązuje dla wszystkich
$Y,W$, to z $aRb_1$ i $aRb_2$ oraz $b_1\ne b_2$
dostajemy sprzeczność, biorąc $Y=\{b_1\}$,
$W=\{b_2\}$: lewy przeciwobraz ich pustego przecięcia
jest pusty, natomiast $a$ należy do prawej strony.
Zatem prawostronna jednoznaczność jest konieczna.
\end{proof}


\begin{przyklad}[Dlaczego przecięcie daje tylko zawieranie?]
Niech $A=\{1,2\}$, $B=\{u\}$ i
$R=\{(1,u),(2,u)\}$. Dla $X=\{1\}$ i $Z=\{2\}$
mamy $X\cap Z=\varnothing$, a więc
$R[X\cap Z]=\varnothing$. Jednak
$R[X]\cap R[Z]=\{u\}$.
Obrazy różnych argumentów mogą się przecinać, nawet gdy
same zbiory argumentów są rozłączne.
\end{przyklad}

\begin{przyklad}[Dlaczego relacyjny przeciwobraz nie zachowuje przecięć]
Niech $A=\{a\}$, $B=\{u,v\}$ oraz
$R=\{(a,u),(a,v)\}\subseteq A\times B$.
Dla $Y=\{u\}$ i $Z=\{v\}$ mamy
$R^{-1}[Y\cap Z]=R^{-1}[\varnothing]=\varnothing$, ale
$R^{-1}[Y]\cap R^{-1}[Z]=\{a\}$.
Powodem jest to, że istnieją \emph{dwa różne} elementy
związane z $a$. Odróżnijmy to od przyszłego pojęcia funkcji,
gdzie każdemu argumentowi przypisany jest tylko jeden wynik.
\end{przyklad}


\begin{przyklad}[Relacja jako lista par]
Niech $A=\{1,2\}$, $B=\{a,b\}$ i
$R=\{(1,a),(2,a)\}\subseteq A\times B$.
Mamy $1Ra$ i $2Ra$, ale $1Rb$ nie zachodzi.
Relacja mówi wyłącznie o tym, które pary należą do $R$;
nie zobowiązuje nas do przypisania każdemu elementowi $A$
jednej i tylko jednej wartości. Ta dodatkowa własność
odróżni funkcje w rozdziale~\ref{ch:functions}.
\end{przyklad}

\section{Relacja odwrotna, identycznościowa i składanie}
\label{sec:relation-algebra}
\begin{definicja}[Relacja odwrotna i identycznościowa]
\label{def:inverse-identity-rel}
Niech $A,B$ będą zbiorami. Dla $R\subseteq A\times B$
relacją odwrotną jest
\[
R^{-1}:=\{(b,a)\in B\times A:(a,b)\in R\}.
\]
Dla dowolnego zbioru $A$ relacja identycznościowa ma postać
$\Delta_A:=\{(a,a):a\in A\}$.
Relacją pełną na $A$ jest $A\times A$, zaś relacją pustą
na $A$ --- $\varnothing\subseteq A\times A$.
\end{definicja}

\begin{twierdzenie}[Obraz względem relacji odwrotnej]
\label{thm:inverse-rel-image}
Niech $A,B$ będą zbiorami, $R\subseteq A\times B$,
$X\subseteq A$ i $Y\subseteq B$. Wówczas
\[
R^{-1}[Y]=\{a\in A:\exists b\in Y\ (a,b)\in R\},
\qquad R[X]=\{b\in B:\exists a\in X\ (b,a)\in R^{-1}\}.
\]
Ponadto $\dom(R)=R^{-1}[B]$ oraz $\ran(R)=R[A]$.
\end{twierdzenie}
\begin{proof}
Weźmy $a\in A$. Z definicji obrazu dla relacji odwrotnej,
$a\in R^{-1}[Y]$ oznacza istnienie $b\in Y$ takiego,
że $(b,a)\in R^{-1}$. Z definicji relacji odwrotnej
zachodzi to wtedy i tylko wtedy, gdy $(a,b)\in R$.
Otrzymaliśmy pierwszą równość. Podobnie $b\in R[X]$
wtedy i tylko wtedy, gdy pewien $a\in X$ spełnia
$(a,b)\in R$, co po odwróceniu pary jest równoważne
$(b,a)\in R^{-1}$. To dowodzi drugiego wzoru.
Jeżeli w pierwszym wzorze podstawimy $Y=B$, otrzymamy
dokładnie definicję $\dom(R)$. Podstawiając $X=A$ w
drugim wzorze, otrzymamy definicję $\ran(R)$.
\end{proof}

\begin{uwaga}[Dwa zastosowania $R^{-1}$]
Symbol $R^{-1}$ oznacza relację powstałą przez odwrócenie
współrzędnych. Zapis $R^{-1}[Y]$ w
Definicji~\ref{def:relation-image} to obraz $Y$ względem
tej właśnie relacji odwrotnej; są to zgodne znaczenia.
\end{uwaga}

\begin{definicja}[Składanie relacji]\label{def:relation-composition}
Niech $A,B,C$ będą zbiorami, $R\subseteq A\times B$
i $S\subseteq B\times C$. Relację
$S\circ R\subseteq A\times C$ definiujemy przez
\[
(a,c)\in S\circ R\quad\Longleftrightarrow\quad
\exists b\in B:\;(a,b)\in R\ \text{oraz}\ (b,c)\in S.
\]
Najpierw przechodzimy relacją $R$ z $A$ do $B$,
a następnie relacją $S$ z $B$ do $C$.
\end{definicja}

\begin{przyklad}[Wyszukiwanie dwustopniowe]
Niech $A=\{a_1,a_2\}$, $B=\{b_1,b_2\}$,
$C=\{c_1,c_2\}$ oraz
$R=\{(a_1,b_1),(a_1,b_2),(a_2,b_2)\}$,
$S=\{(b_1,c_1),(b_2,c_2)\}$.
Wtedy
$S\circ R=\{(a_1,c_1),(a_1,c_2),(a_2,c_2)\}$.
Na przykład $(a_1,c_2)$ należy do składania, gdyż
możemy wskazać pośredni element $b_2$.
\end{przyklad}

\begin{center}
\begin{tikzpicture}[>=Stealth, every node/.style={font=\small}]
\node (a1) at (0,1.2) {$a_1$}; \node (a2) at (0,0) {$a_2$};
\node (b1) at (3,1.2) {$b_1$}; \node (b2) at (3,0) {$b_2$};
\node (c1) at (6,1.2) {$c_1$}; \node (c2) at (6,0) {$c_2$};
\draw[->,thick] (a1)--(b1); \draw[->,thick] (a1)--(b2);
\draw[->,thick] (a2)--(b2); \draw[->,thick] (b1)--(c1);
\draw[->,thick] (b2)--(c2);
\node at (1.5,1.7) {$R$}; \node at (4.5,1.7) {$S$};
\end{tikzpicture}

{\small Przejścia przez punkt pośredni opisują relację $S\circ R$.}
\end{center}

\begin{lemat}[Monotoniczność składania]
\label{lem:composition-monotone}
Niech $A,B,C$ będą zbiorami, $R_1,R_2\subseteq A\times B$
oraz $S_1,S_2\subseteq B\times C$. Jeśli
$R_1\subseteq R_2$ i $S_1\subseteq S_2$,
to $S_1\circ R_1\subseteq S_2\circ R_2$.
\end{lemat}
\begin{proof}
Weźmy $(a,c)\in S_1\circ R_1$. Istnieje wtedy $b\in B$
takie, że $(a,b)\in R_1$ i $(b,c)\in S_1$.
Z założonych zawierań wynika $(a,b)\in R_2$
oraz $(b,c)\in S_2$. Ten sam element pośredni $b$
świadczy zatem o $(a,c)\in S_2\circ R_2$.
\end{proof}

\begin{twierdzenie}[Łączność składania relacji]\label{thm:relation-assoc}
Niech $A,B,C,D$ będą zbiorami, $R\subseteq A\times B$,
$S\subseteq B\times C$ i $T\subseteq C\times D$. Wtedy
\[
T\circ(S\circ R)=(T\circ S)\circ R.
\]
\end{twierdzenie}
\begin{proof}
Ustalmy $(a,d)\in A\times D$.
Z definicji składania $(a,d)\in T\circ(S\circ R)$
wtedy i tylko wtedy, gdy istnieje $c\in C$ taki, że
$(a,c)\in S\circ R$ oraz $(c,d)\in T$.
Pierwszy z tych warunków oznacza istnienie $b\in B$
z $(a,b)\in R$ i $(b,c)\in S$. Łącznie otrzymujemy:
istnieją $b\in B$ i $c\in C$ takie, że
\[
(a,b)\in R,\qquad (b,c)\in S,\qquad (c,d)\in T.
\]
To samo zdanie możemy odczytać inaczej: istnieje $b\in B$
taki, że $(a,b)\in R$ oraz istnieje $c\in C$,
dla którego $(b,c)\in S$ i $(c,d)\in T$.
Druga część jest równoważna $(b,d)\in T\circ S$.
Stąd całość jest równoważna
$(a,d)\in(T\circ S)\circ R$.
Ponieważ każdy krok można odwrócić, obie relacje zawierają
te same pary.
\end{proof}

\begin{twierdzenie}[Relacje identycznościowe i odwracanie składania]
\label{thm:relation-identities}
Niech $A,B$ będą zbiorami i $R\subseteq A\times B$.
Wówczas $R\circ\Delta_A=R=\Delta_B\circ R$ oraz
$(R^{-1})^{-1}=R$.
Jeżeli dodatkowo $S\subseteq B\times C$, to
\[
(S\circ R)^{-1}=R^{-1}\circ S^{-1}.
\]
\end{twierdzenie}
\begin{proof}
Dla $(a,b)\in A\times B$ warunek
$(a,b)\in R\circ\Delta_A$ oznacza, że istnieje
$a'\in A$ z $(a,a')\in\Delta_A$ i $(a',b)\in R$.
Pierwszy warunek wymusza $a'=a$, a drugi redukuje się do
$(a,b)\in R$. Dowodzi to pierwszej równości.
Analogicznie przy $\Delta_B\circ R$ pośredni element
$b'$ musi być równy $b$. Dwukrotne odwrócenie pary
$(a,b)$ przywraca tę samą parę, więc $(R^{-1})^{-1}=R$.

Wreszcie $(c,a)\in(S\circ R)^{-1}$
wtedy i tylko wtedy, gdy $(a,c)\in S\circ R$,
czyli gdy dla pewnego $b\in B$ zachodzą
$(a,b)\in R$ i $(b,c)\in S$.
Po odwróceniu współrzędnych są to warunki
$(c,b)\in S^{-1}$ i $(b,a)\in R^{-1}$,
a więc $(c,a)\in R^{-1}\circ S^{-1}$.
Równoważność działa w obu kierunkach.
\end{proof}

\section{Własności relacji na zbiorze}\label{sec:relation-properties}
\begin{definicja}[Sześć podstawowych własności]\label{def:relation-properties}
Niech $A$ będzie zbiorem, a $R\subseteq A\times A$.
Relację $R$ nazywamy:
\begin{enumerate}[label=(\roman*)]
\item \emph{zwrotną}, gdy $\forall a\in A\;aRa$;
\item \emph{przeciwzwrotną}, gdy $\forall a\in A\;\neg(aRa)$;
\item \emph{symetryczną}, gdy $\forall a,b\in A\;(aRb\Rightarrow bRa)$;
\item \emph{antysymetryczną}, gdy
$\forall a,b\in A\;((aRb\land bRa)\Rightarrow a=b)$;
\item \emph{asymetryczną}, gdy
$\forall a,b\in A\;(aRb\Rightarrow\neg(bRa))$;
\item \emph{przechodnią}, gdy
$\forall a,b,c\in A\;((aRb\land bRc)\Rightarrow aRc)$.
\end{enumerate}
\end{definicja}

\begin{uwaga}[Antysymetria nie znaczy braku symetrii]
Relacja $\le$ na $\R$ jest antysymetryczna, ponieważ
$x\le y$ i $y\le x$ implikują $x=y$. Jest też zwrotna:
dla każdej liczby $x$ zachodzi $x\le x$. W szczególności,
gdy $x=y$, oba warunki z definicji antysymetrii są spełnione.
Antysymetria dopuszcza zatem pary $(x,x)$, natomiast
asymetria wyklucza je całkowicie.
\end{uwaga}

\begin{twierdzenie}[Asymetria i podstawowe własności]
\label{thm:asymmetric}
Niech $A$ będzie zbiorem, a $R\subseteq A\times A$.
Wówczas $R$ jest asymetryczna wtedy i tylko wtedy, gdy
jest jednocześnie przeciwzwrotna i antysymetryczna.
Ponadto $R$ jest symetryczna wtedy i tylko wtedy,
gdy $R=R^{-1}$, a przechodnia wtedy i tylko wtedy,
gdy $R\circ R\subseteq R$.
\end{twierdzenie}
\begin{proof}
Załóżmy asymetrię. Gdyby istniał $a\in A$ z $aRa$,
to podstawiając $b=a$ do definicji asymetrii otrzymalibyśmy
$\neg(aRa)$, sprzeczność. Stąd przeciwzwrotność.
Jeśli $aRb$ i $bRa$, to z pierwszego warunku i asymetrii
wynika $\neg(bRa)$, także sprzeczność; sytuacja ta
nigdy więc nie występuje. W szczególności implikacja
$(aRb\land bRa)\Rightarrow a=b$ jest zawsze prawdziwa,
czyli $R$ jest antysymetryczna.

Odwrotnie, załóżmy przeciwzwrotność i antysymetrię.
Weźmy $a,b\in A$ z $aRb$. Jeśli również $bRa$,
antysymetria dałaby $a=b$, a wówczas $aRa$,
co przeczy przeciwzwrotności. Zatem $\neg(bRa)$,
czyli $R$ jest asymetryczna.

Dla symetrii zauważmy, że $R\subseteq R^{-1}$ oznacza,
że z $(a,b)\in R$ wynika $(a,b)\in R^{-1}$,
czyli $(b,a)\in R$. To dokładnie definicja symetrii.
Jeżeli $R\subseteq R^{-1}$, odwrócenie obu stron
daje $R^{-1}\subseteq R$; otrzymujemy więc $R=R^{-1}$.
Odwrotnie, jeżeli $R=R^{-1}$, to w szczególności
$R\subseteq R^{-1}$, a wyżej pokazaliśmy, że takie zawieranie
jest równoważne symetrii relacji $R$.

Wreszcie para $(a,c)$ należy do $R\circ R$ dokładnie
wtedy, gdy istnieje $b\in A$ z $aRb$ i $bRc$.
Warunek $R\circ R\subseteq R$ mówi zatem, że dla
każdej takiej pary zachodzi $aRc$, co jest definicją
przechodniości.
\end{proof}

\begin{przyklad}[Jedna relacja, wiele własności]
Na $\Z$ relacja $<$ jest przeciwzwrotna, asymetryczna
i przechodnia; nie jest zwrotna ani symetryczna.
Relacja $\le$ jest zwrotna, antysymetryczna i przechodnia.
Relacja „mieć tę samą parzystość” na $\Z$ jest zwrotna,
symetryczna i przechodnia, ale nie jest antysymetryczna,
gdyż $2$ i $4$ są różne, a mają tę samą parzystość.
Na $\Z$ relacja podzielności nie jest antysymetryczna,
gdyż $2\mid(-2)$ i $(-2)\mid2$, ale $2\ne-2$.
\end{przyklad}

\begin{przyklad}[Niezależność przechodniości od innych własności]
Niech $A=\{1,2,3\}$, $R=\{(1,2),(2,3)\}$.
Relacja ta jest asymetryczna, a więc przeciwzwrotna
i antysymetryczna, lecz nie jest przechodnia:
$1R2$ i $2R3$, podczas gdy $1R3$ nie zachodzi.
Relacja pusta jest z kolei symetryczna, asymetryczna
i przechodnia na dowolnym $A$; gdy $A\ne\varnothing$,
nie jest zwrotna. Na $A=\varnothing$ staje się również
zwrotna, ponieważ nie istnieje element, który naruszałby
warunek zwrotności.
\end{przyklad}

\section{Najmniejsze relacje o zadanych własnościach}
\label{sec:relation-closures}
Często początkowo znamy tylko bezpośrednie połączenia między
obiektami. Chcemy jednak rozszerzyć je tak, aby wolno było
traktować połączenia jako osiągalność po dowolnej liczbie kroków.
To podstawowa motywacja domknięcia przechodniego.

\begin{definicja}[Potęgi relacji i ścieżki]
\label{def:relation-powers}
Niech $A$ będzie zbiorem i $R\subseteq A\times A$.
Definiujemy rekurencyjnie
$R^0:=\Delta_A$ i $R^{n+1}:=R\circ R^n$ dla
$n\in\N:=\{0,1,2,\ldots\}$. Para $(a,b)\in R^n$
oznacza istnienie ścieżki o długości $n$ od $a$ do $b$,
tj. elementów $a=a_0,a_1,\ldots,a_n=b$ takich,
że $a_{j-1}Ra_j$ dla $j=1,\ldots,n$.
Dla $n=0$ ścieżka składa się z jednego punktu $a=b$.
\end{definicja}

\begin{lemat}[Ścieżki a potęgi relacji]\label{lem:relation-paths-xlii}
Niech $A$ będzie zbiorem, $R\subseteq A\times A$,
$a,b\in A$ oraz $n\in\N$. Wtedy $(a,b)\in R^n$
wtedy i tylko wtedy, gdy istnieją
$a_0,\ldots,a_n\in A$ takie, że
$a_0=a$, $a_n=b$ oraz $a_{j-1}Ra_j$ dla
każdego $j\in\{1,\ldots,n\}$.
Dla $n=0$ ostatnia lista warunków jest pusta,
ale nadal wymagamy $a_0=a=b$.
\end{lemat}
\begin{proof}
Dowodzimy równoważności przez indukcję po $n$.
Dla $n=0$ z definicji $R^0=\Delta_A$ mamy
$(a,b)\in R^0$ wtedy i tylko wtedy, gdy $a=b$.
To dokładnie warunek istnienia ścieżki długości zero,
składającej się z jednego wierzchołka $a_0=a=b$.

Załóżmy twierdzenie dla ustalonej liczby $n$.
Z definicji $R^{n+1}=R\circ R^n$ para $(a,b)$
należy do $R^{n+1}$ wtedy i tylko wtedy, gdy
istnieje punkt $c\in A$ taki, że $(a,c)\in R^n$
i $cRb$. Z założenia indukcyjnego pierwszy warunek
jest równoważny istnieniu ciągu
$a=a_0,\ldots,a_n=c$ z kolejnymi krawędziami w $R$.
Dołączenie $a_{n+1}=b$ daje ścieżkę długości $n+1$
od $a$ do $b$. Odwrotnie, ze ścieżki
$a=a_0,\ldots,a_n,a_{n+1}=b$ wybieramy $c=a_n$;
jej pierwsze $n$ krawędzi, na mocy założenia
indukcyjnego, świadczy o $(a,c)\in R^n$,
a ostatnia o $cRb$. Definicja składania daje
$(a,b)\in R^{n+1}$. To zamyka obie implikacje.
\end{proof}


\begin{lemat}[Potęgi relacji i łączenie ścieżek]
\label{lem:relation-powers}
Niech $A$ będzie zbiorem, $R\subseteq A\times A$
i $m,n\in\N$. Wówczas zachodzi $R^{m+n}=R^m\circ R^n$.
\end{lemat}
\begin{proof}
Prowadzimy indukcję po $m$. Dla $m=0$ prawa strona
jest równa $\Delta_A\circ R^n=R^n$ na mocy
Twierdzenia~\ref{thm:relation-identities}; lewa strona
to $R^{0+n}=R^n$. Załóżmy równość dla ustalonego $m$
i każdego $n$. Wówczas, korzystając z definicji potęgi
i łączności składania
(Twierdzenie~\ref{thm:relation-assoc}), mamy
\[
R^{(m+1)+n}=R\circ R^{m+n}
=R\circ(R^m\circ R^n)
=(R\circ R^m)\circ R^n=R^{m+1}\circ R^n.
\]
Krok indukcyjny kończy dowód.
\end{proof}

\begin{definicja}[Domknięcia zwrotne, symetryczne, przechodnie]
\label{def:relation-closure}
Niech $A$ będzie zbiorem i $R\subseteq A\times A$.
Definiujemy
\[
R^{\mathrm{ref}}:=R\cup\Delta_A,\qquad
R^{\mathrm{sym}}:=R\cup R^{-1},\qquad
R^+:=\bigcup_{n\ge1}R^n,\qquad
R^*:=\bigcup_{n\ge0}R^n.
\]
Relacje $R^+$ oraz $R^*$ nazywamy odpowiednio
\emph{domknięciem przechodnim} i
\emph{domknięciem zwrotno-przechodnim} relacji $R$.
\end{definicja}

\begin{twierdzenie}[Minimalność domknięć]\label{thm:relation-closure}
Niech $A$ będzie zbiorem i $R\subseteq A\times A$.
Wtedy $R^{\mathrm{ref}}$ jest najmniejszą (względem zawierania)
relacją zwrotną na $A$, która zawiera $R$,
a $R^{\mathrm{sym}}$ jest najmniejszą relacją symetryczną
na $A$, która zawiera $R$. Ponadto $R^+$ jest najmniejszą
relacją przechodnią na $A$ zawierającą $R$,
a $R^*$ jest najmniejszą relacją zwrotną i przechodnią
na $A$ zawierającą $R$.
\end{twierdzenie}
\begin{proof}
Relacja $R\cup\Delta_A$ zawiera $R$ i wszystkie pary
$(a,a)$, więc jest zwrotna. Jeżeli zwrotna relacja
$S\subseteq A\times A$ zawiera $R$, musi zawierać
każdą parę przekątnej $\Delta_A$, a więc również
$R\cup\Delta_A$. To dowodzi minimalności
$R^{\mathrm{ref}}$.

Jeśli $(a,b)\in R\cup R^{-1}$, to
$(a,b)\in R$ lub $(b,a)\in R$. W pierwszym przypadku
$(b,a)\in R^{-1}$, a w drugim $(b,a)\in R$.
Zatem $R\cup R^{-1}$ jest symetryczna. Jeżeli symetryczna
relacja $S$ zawiera $R$, to zawiera także $R^{-1}$,
więc zawiera sumę $R\cup R^{-1}$.

Przejdźmy do $R^+$. Ponieważ $R^1=R$, mamy
$R\subseteq R^+$. Jeśli $(a,b)\in R^+$ i $(b,c)\in R^+$,
wybierzmy dodatnie liczby naturalne $m,n$ takie, że
$(a,b)\in R^m$ i $(b,c)\in R^n$.
Z definicji składania $(a,c)\in R^n\circ R^m$.
Na mocy Lematu~\ref{lem:relation-powers} prawa strona
jest równa $R^{m+n}$. Ponieważ $m+n\ge1$,
mamy $(a,c)\in R^+$. Zatem $R^+$ jest przechodnia.
Niech teraz $S$ będzie dowolną przechodnią relacją
zawierającą $R$. Pokażemy indukcyjnie, że
$R^n\subseteq S$ dla wszystkich $n\ge1$.
Dla $n=1$ jest to założenie $R\subseteq S$.
Jeśli $R^n\subseteq S$, to
$R^{n+1}=R\circ R^n\subseteq S\circ S\subseteq S$:
pierwsze zawieranie wynika z Lematu~\ref{lem:composition-monotone},
a drugie z przechodniości $S$ (zob.
Twierdzenie~\ref{thm:asymmetric}).
Stąd każdy składnik sumy $R^+$ zawiera się w $S$,
więc $R^+\subseteq S$.

Wreszcie $R^*=\Delta_A\cup R^+$.
Jest zwrotna, gdyż zawiera $\Delta_A$.
Jest też przechodnia: przy złożeniu dwóch par z
$R^*$ możemy dobrać $m,n\ge0$ i użyć
$R^n\circ R^m=R^{m+n}\subseteq R^*$.
Jeśli $S$ jest zwrotna i przechodnia oraz zawiera $R$,
to zawiera $R^+$ z poprzedniego argumentu i
$\Delta_A$ ze zwrotności, a zatem
$R^*=R^+\cup\Delta_A\subseteq S$.
\end{proof}

\begin{przyklad}[Zależność w grafie]
Niech $A=\{a,b,c,d\}$, $R=\{(a,b),(b,c),(c,d)\}$.
W $R^+$ pojawią się między innymi pary $(a,c)$,
$(b,d)$ oraz $(a,d)$, których nie ma w $R$.
Oznacza to, że możemy przejść od $a$ do $d$ po trzech
krawędziach. W $R^*$ dodatkowo obecne będą
$(a,a),(b,b),(c,c),(d,d)$, odpowiadające ścieżkom
o długości zero.
\end{przyklad}

\begin{uwaga}[Zapowiedź relacji równoważności i porządków]
Relacja jest \emph{relacją równoważności}, gdy jest
zwrotna, symetryczna i przechodnia; jej klasy i ilorazy
zbadamy w rozdziale 5. Relacja jest \emph{porządkiem
częściowym}, gdy jest zwrotna, antysymetryczna i przechodnia;
porządki, kresy i diagramy Hassego będą tematem rozdziału 6.
Niniejszy rozdział dostarcza narzędzi potrzebnych do obu teorii.
\end{uwaga}

\begin{twierdzenie}[Najmniejsza relacja równoważności generowana przez relację]
\label{thm:equiv-generated-prelim}
Niech $A$ będzie zbiorem, a $R\subseteq A\times A$ relacją.
Wtedy relacja $E:=(R\cup R^{-1})^*$ jest zwrotna, symetryczna
i przechodnia. Jest też najmniejszą relacją o tych trzech
własnościach zawierającą $R$.
\end{twierdzenie}
\begin{proof}
Oznaczmy $S=R\cup R^{-1}$. Z definicji $S$ jest symetryczna:
jeśli $(a,b)\in S$, to $(b,a)\in S$. Relacja $S^*$ zawiera $S$
i jest zwrotna i przechodnia z
Twierdzenia~\ref{thm:relation-closure}. Aby wykazać symetrię $S^*$,
weźmy $(a,b)\in S^*$. Istnieje wtedy ciąg
$a=a_0,a_1,\ldots,a_n=b$, w którym każda kolejna para
$(a_{j-1},a_j)$ należy do $S$ (dla $n=0$ mamy $a=b$).
Dzięki symetrii $S$ wszystkie pary
$(a_j,a_{j-1})$ też należą do $S$. Odwracając kolejność
punktów, otrzymujemy ścieżkę od $b$ do $a$, zatem
$(b,a)\in S^*$. To dowodzi symetrii.

Niech teraz $T$ będzie dowolną relacją zwrotną, symetryczną
i przechodnią zawierającą $R$. Ze zwrotności i symetrii wynika
$R^{-1}\subseteq T$, zatem $S\subseteq T$.
Z minimalności domknięcia zwrotno-przechodniego
(\ref{thm:relation-closure}) dostajemy $S^*\subseteq T$.
Dokładnie to znaczy, że $E$ jest najmniejszą z żądanych relacji.
\end{proof}

\begin{przyklad}[Spójne składowe grafu]
Dla $A=\{1,2,3,4,5\}$ i $R=\{(1,2),(2,3),(4,5)\}$
relacja $(R\cup R^{-1})^*$ łączy każdą parę elementów z
$\{1,2,3\}$ oraz każdą parę z $\{4,5\}$, ale nie łączy
elementów należących do różnych tych grup.
W rozdziale o relacjach równoważności pokażemy,
że te dwie grupy są klasami równoważności.
\end{przyklad}

\begin{center}
\begin{tikzpicture}[>=Stealth,every node/.style={font=\small}]
  \node[circle,draw,minimum size=6.5mm] (v1) at (0,0) {$1$};
  \node[circle,draw,minimum size=6.5mm] (v2) at (1.5,.65) {$2$};
  \node[circle,draw,minimum size=6.5mm] (v3) at (3,0) {$3$};
  \node[circle,draw,minimum size=6.5mm] (v4) at (5.0,.65) {$4$};
  \node[circle,draw,minimum size=6.5mm] (v5) at (6.5,0) {$5$};
  \draw[very thick] (v1)--(v2)--(v3);
  \draw[very thick] (v4)--(v5);
  \node[below=8mm of v2] {$\{1,2,3\}$};
  \node at (5.75,-.9) {$\{4,5\}$};
  \node[font=\small,anchor=north] at (3.3,-1.6)
    {Dwie składowe relacji $\bigl(R\cup R^{-1}\bigr)^*$.};
\end{tikzpicture}
\end{center}


\section{Relacje a macierze zero-jedynkowe}
\label{sec:relation-matrices}
\begin{przyklad}[Konkretne obliczenie składania]
Dla skończonych zbiorów $A=\{a_1,\ldots,a_m\}$ i
$B=\{b_1,\ldots,b_n\}$ relację $R\subseteq A\times B$
można zakodować macierzą $M_R\in\{0,1\}^{m\times n}$:
$(M_R)_{ij}=1$ wtedy i tylko wtedy, gdy $a_iRb_j$.
Jeśli $S\subseteq B\times C$, to macierz składania
ma współczynniki
\[
(M_{S\circ R})_{ik}
=\bigvee_{j=1}^{n}\bigl((M_R)_{ij}\land(M_S)_{jk}\bigr),
\]
gdzie $\lor$ i $\land$ oznaczają operacje logiczne.
Przy $B=\varnothing$ mamy $n=0$ i przyjmujemy,
że alternatywa pustej rodziny ma wartość $0$;
odpowiada to brakowi elementu pośredniego i daje
relację pustą $S\circ R=\varnothing$.
Aby współczynnik był równy $1$, musi istnieć przynajmniej
jedno $j$ takie, że oba czynniki mają wartość $1$;
jest to dokładnie warunek istnienia punktu pośredniego
$b_j$ w definicji składania.

Na przykład dla $A=B=C=\{1,2\}$ i
$R=\{(1,1),(1,2)\}$ oraz $S=\{(1,2),(2,2)\}$ mamy
\[
M_R=\begin{pmatrix}1&1\\0&0\end{pmatrix},\qquad
M_S=\begin{pmatrix}0&1\\0&1\end{pmatrix},\qquad
M_{S\circ R}=\begin{pmatrix}0&1\\0&0\end{pmatrix}.
\]
Takie macierze i iloczyn boolowski stanowią podstawę
algorytmów obliczania osiągalności w skończonych grafach.
\end{przyklad}

\begin{twierdzenie}[Macierz relacji odwrotnej]\label{thm:lxiv-relation-matrix-transpose}
Niech $A=\{a_1,\ldots,a_m\}$ i $B=\{b_1,\ldots,b_n\}$
będą skończonymi zbiorami z ustalonymi wyliczeniami,
a $R\subseteq A\times B$ relacją. Wówczas macierz relacji
odwrotnej $R^{-1}\subseteq B\times A$, przy tych samych
wyliczeniach, jest transpozycją macierzy $R$:
$M_{R^{-1}}=M_R^{\mathsf T}$. Równość obowiązuje także
wtedy, gdy jeden z tych zbiorów jest pusty.
\end{twierdzenie}
\begin{proof}
Dla dowolnych $1\le j\le n$ i $1\le i\le m$,
z definicji macierzy relacji mamy
$(M_{R^{-1}})_{ji}=1$ wtedy i tylko wtedy, gdy
$b_jR^{-1}a_i$. Z definicji relacji odwrotnej jest to
równoważne $a_iRb_j$, a więc $(M_R)_{ij}=1$.
Ponieważ oba współczynniki należą do $\{0,1\}$,
są równe. Jeżeli $m=0$ lub $n=0$, nie istnieje
żadna para indeksów do sprawdzenia, a obie macierze
mają odpowiednio rozmiary $n\times m$ i $m\times n$;
transpozycja nadal ma wymagany rozmiar i wpisy.
\end{proof}
\begin{przyklad}[Transpozycja a kierunek strzałek]\label{ex:lxiv-relation-matrix-transpose}
Niech $A=\{1,2\}$, $B=\{u,v,w\}$ i
$R=\{(1,u),(1,w),(2,v)\}$. Wówczas
\[
M_R=\begin{pmatrix}1&0&1\\0&1&0\end{pmatrix},\qquad
M_{R^{-1}}=\begin{pmatrix}1&0\\0&1\\1&0\end{pmatrix}.
\]
Na przykład wpis $(3,1)$ drugiej macierzy wynosi $1$,
ponieważ $wR^{-1}1$, czyli $(1,w)\in R$.
\end{przyklad}

\begin{uwaga}[Podsumowanie metody dowodzenia]
Najczęstsza technika w tych rozdziałach to wzięcie
dowolnego elementu $x$ lub pary $(a,b)$ i zastąpienie
przynależności do zbioru rozwinięciem jego definicji.
W szczególności prawa składania relacji są prawami
manipulowania kwantyfikatorami egzystencjalnymi:
„istnieje element pośredni”. W następnych rozdziałach
będziemy coraz częściej korzystać z tej metody.
\end{uwaga}
